\chapter{The two heads, and the loop that trains them}
\label{ch:heads}

\begin{objectives}
\begin{itemize}
\item Describe how a $64\times768$ block of numbers becomes a probability over $\sim\!1858$
  moves and a win/draw/loss triple.
\item Verify, on real output, that the policy is a proper distribution over legal moves.
\item Explain how the same network can be right in one head and wrong in the other about the
  same move --- with a real example.
\item State the self-play training loop as three equations and explain why it improves.
\item Say precisely whose chess the policy head models, and what follows for a human training
  tool.
\end{itemize}
\end{objectives}

\section{Two questions, one body}

Chapter~\ref{ch:problem} said every engine must answer two questions: how good is this, and
where should I look? Leela answers both with one network and two small \emph{heads} attached
to the top of it.

\begin{center}
\begin{tikzpicture}[
  box/.style={draw, rounded corners=2pt, minimum height=7mm, align=center, font=\small},
  arr/.style={-{Latex[length=2mm]}, thick}]
\node[box, fill=black!4, text width=32mm] (in) at (0,0) {board $+$ history\\\footnotesize $64\times768$};
\node[box, fill=BuildBlue!12, text width=38mm, minimum height=16mm] (body) at (4.4,0)
  {\textbf{15 encoder layers}\\\footnotesize attention + feed-forward};
\node[box, fill=IdeaGreen!12, text width=34mm] (pol) at (9.4,0.95)
  {\textbf{policy head}\\\footnotesize $P(a\given s)$, $\sim$1858 slots};
\node[box, fill=DataTeal!12, text width=34mm] (val) at (9.4,-0.35)
  {\textbf{value head}\\\footnotesize $(w, d, l)$};
\node[box, fill=black!4, text width=34mm] (mlh) at (9.4,-1.6)
  {\footnotesize moves-left head: $M$};
\draw[arr] (in) -- (body);
\draw[arr] (body.east) -- (pol.west);
\draw[arr] (body.east) -- (val.west);
\draw[arr] (body.east) -- (mlh.west);
\end{tikzpicture}
\end{center}

Everything below the heads is shared. That is not a saving; it is the point. The
representation that lets you say "White is winning" is largely the representation that lets
you say "and the move is Qb8+", so the two tasks teach each other. This is called multi-task
learning, and in AlphaZero-style training it is essential rather than incidental.

\section{The policy head}

\subsection{Naming a move with a number}

A probability distribution needs a fixed set of outcomes, so every possible chess move needs
its own slot. Leela uses a fixed vocabulary of \textbf{1858} move slots, covering every
(from-square, to-square) pair that any piece could ever traverse, plus underpromotions.

Most are illegal in any given position. So the head produces a raw score --- a
\textbf{logit} --- for all 1858, and then:
\begin{enumerate}
\item set the logits of illegal moves to $-\infty$;
\item softmax over what remains.
\end{enumerate}
\begin{equation}
  \label{eq:policyhead}
  P(a \given s) \;=\;
  \frac{\exp\bigl(z_a\bigr)}{\sum_{b \in \mathcal{A}(s)}\exp\bigl(z_b\bigr)},
  \qquad z = \text{(linear map of }X^{(15)}).
\end{equation}
The masking matters: the network is never asked to learn the rules of chess, only which legal
move to prefer. This is why the priors you saw in Chapter~\ref{ch:puct} always add up:

\begin{realdata}[title={Real engine output: the priors are a proper distribution}]
Summing the reported $P$ over all legal moves in each position of this book's data set:
\begin{center}
\begin{tabular}{@{}lrr@{}}
\toprule
Position & legal moves & $\sum_a P(a\given s)$ \\
\midrule
Initial position & 20 & $100.00\%$ \\
K+P endgame & 4 & $100.00\%$ \\
Drawn opposition ending & 5 & $99.99\%$ \\
Morphy, before 16.Qb8+ & 46 & $100.00\%$ \\
Morphy, before 12.O-O-O & 49 & $99.96\%$ \\
\bottomrule
\end{tabular}
\end{center}
The shortfalls are rounding in the printed two-decimal percentages, not a modelling gap. Note
also that every legal move appears with a non-zero prior --- there is no hard filter anywhere
(Chapter~\ref{ch:puct}, Attempt 1).
\end{realdata}

\subsection{What the policy is a model of}

Read Equation~\eqref{eq:policyhead} together with what the training target will turn out to be
(§12.4) and you get a statement that is easy to misremember:

\begin{keyidea}
The policy head predicts \textbf{which move a full search by this same engine would end up
preferring}. It is a fast approximation of its own slow self. It is \emph{not} a model of
human choice, \emph{not} a goodness score, and \emph{not} calibrated to any rating band.
\end{keyidea}

That has a direct consequence for your project which is worth stating plainly, because it
cuts against the natural reading of a "policy blindness" metric. When your tool says "you
played a move Leela's instinct rated at $0.4\%$", the honest gloss is \emph{"a move this
engine would rarely choose"}, not \emph{"a move a strong human would rarely choose"}. Those
overlap heavily --- the network is very strong, and strong play is strong play --- but they
are not the same claim, and the gap is largest precisely in the sharp, unusual positions the
project cares about most.

\section{The value head}

The value head is smaller: a couple of linear layers, then a three-way softmax:
\begin{equation}
  \label{eq:valuehead}
  (w, d, l) \;=\; \softmax\bigl(W_2\,\relu(W_1 \bar{\vv{x}} + \vv{b}_1) + \vv{b}_2\bigr),
  \qquad V = w - l ,
\end{equation}
where $\bar{\vv{x}}$ is a pooled summary of the 64 square-vectors. Three outputs, softmaxed,
so they are non-negative and sum to one --- exactly the WDL triple of
Chapter~\ref{ch:currency}.

Older networks emitted a single scalar $V$ instead. Splitting it into three was a real gain,
and not only for display: it lets the network express \emph{how} it expects a position to
finish, and it makes contempt and draw-avoidance expressible at search time without
retraining.

\subsection{The moves-left head}

A third head predicts the number of plies remaining, the \texttt{M} column you have been
ignoring:

\begin{realdata}[title={Real engine output: the moves-left head}]
\begin{center}
\begin{tabular}{@{}lrl@{}}
\toprule
Position & $M$ (plies) & \\
\midrule
Initial position & $168.0$ & about 84 moves --- a full game \\
Morphy, before 12.O-O-O & $71.0$ & \\
Morphy, before 16.Qb8+ & $86.6$ & the net does not see the mate \\
K+P endgame (won) & $22.5$ & \\
Drawn opposition ending & $19.2$ & \\
\bottomrule
\end{tabular}
\end{center}
\end{realdata}

Its purpose is practical: among equally winning moves, prefer the one that finishes sooner
(and among lost ones, the one that drags on). Without it, an engine in a won position happily
shuffles, because every shuffle is also winning. It is a small head with an outsized effect
on whether an engine's play looks purposeful.

Notice the third row: in a position with mate in two, the moves-left head predicts $86.6$
plies. The network does not see the mate --- consistent with its $V$ of $+0.62$
(Chapter~\ref{ch:currency}) and its $1.6\%$ policy for the mating move. All three heads agree
in their blindness, and search corrected all three at once.

\section{One network, two opinions: how the heads disagree}

The heads share a body but have separate output layers and separate training targets, so they
can disagree. Chapter~\ref{ch:byhand} contains a clean example:

\begin{realdata}[title={Real engine output: the value head knows what the policy head does not}]
K+P endgame, the two moves that throw away the win:
\begin{center}
\begin{tabular}{@{}lrrrl@{}}
\toprule
Move & $P$ (policy) & $V$ of the resulting position & $d$ & Stockfish d30 \\
\midrule
Kd5 & $5.26\%$ & $0.000$ & $1.000$ & $0.00$ (draw) \\
Kf5 & $5.38\%$ & $0.000$ & $0.999$ & $0.00$ (draw) \\
\bottomrule
\end{tabular}
\end{center}
The value head is certain --- draw probability $1.000$, evaluation exactly zero, and Stockfish
agrees. The policy head still assigns these moves a tenth of its attention.
\end{realdata}

Why? Because they are answering different questions. The value head is asked \emph{how does
this position finish}, and a king-and-pawn structure with the defender holding the opposition
is a shape it has seen a hundred million times. The policy head is asked \emph{which move
would search choose}, at a moment when it has not looked at any resulting position at all ---
it must judge Kd5 without evaluating what Kd5 leads to.

\begin{keyidea}
\textbf{The policy sees the move; the value sees the consequence.} A single visit is exactly
the operation that lets the second inform the first --- which is why in
Chapter~\ref{ch:byhand} two visits out of 800 were enough to eliminate both errors. And it is
why, in the position where \emph{both} heads are wrong (the Morphy mate), the correction took
188 nodes instead of 2.
\end{keyidea}

\section{The training loop}

Now the part that makes the whole thing self-improving. Leela learns from games it plays
against itself, with no human games and no human evaluations.

\subsection{Where the training targets come from}

Play a self-play game. At every position $s_t$:

\begin{itemize}
\item run a search of a few hundred nodes;
\item record the \textbf{visit distribution} $\pi_t(a) = N(s_t,a)/N(s_t)$ ---
  Chapter~\ref{ch:engineering}'s Equation~\eqref{eq:temperature} with $\tau = 1$;
\item play a move sampled from $\pi_t$;
\item at the end of the game, record the result $z \in \{+1, 0, -1\}$ from each position's
  point of view.
\end{itemize}

Then train the network to make its raw outputs match:
\begin{equation}
  \label{eq:azloss}
  L \;=\;
  \underbrace{-\sum_a \pi_t(a)\ln P(a\given s_t)}_{\text{policy: cross-entropy against the visits}}
  \;+\;
  \underbrace{\bigl(V(s_t) - z\bigr)^2}_{\text{value: match the result}}
  \;+\;
  \underbrace{\lambda\lVert\theta\rVert^2}_{\text{regularisation}} .
\end{equation}

\subsection{Why this improves anything}

This is the part that looks circular and is not. The key observation:

\begin{established}
\textbf{Search is a policy improvement operator.} The visit distribution $\pi$ produced by a
search of $N$ nodes is a \emph{better} policy than the raw $P$ that guided it --- because $\pi$
is $P$ reshaped by actual evaluations of actual resulting positions
(Chapter~\ref{ch:puct} §6.3: visits reproduce the prior only when search learns nothing).
Training $P$ towards $\pi$ therefore moves the network towards a policy it did not previously
have. The improved $P$ then makes the next search better, which makes the next $\pi$ better
still.
\end{established}

You have already seen this operator at work, with numbers. In the initial position the raw
policy said e4 $22.2\%$; after 1600 nodes the visits said $35.7\%$. The training target for
that position would be $35.7\%$ --- and a network trained on many such examples drifts towards
the opinion its own search keeps arriving at. In the K+P ending the policy said $10.6\%$ for
the two drawing moves; the visits said $0.4\%$. Feed that back and the policy stops liking
them.

\begin{keyidea}
\textbf{The loop is: intuition guides calculation; calculation corrects intuition.} That is
not merely a nice description of AlphaZero --- it is a fair description of what a chess player
does over years, and it is exactly what your project is trying to do for a human. Your
"intuition drill" (guess Leela's top policy move) is the human-facing version of
Equation~\eqref{eq:azloss}: it exposes where the student's $P$ differs from a good $\pi$.
\end{keyidea}

\subsection{Honest caveats}

\begin{researchfinding}
The AlphaZero result (Silver et al., 2017--18) is real and reproduced many times, including by
the Leela project in the open. But note the conditions: the loop needs an enormous number of
self-play games to work, the improvement is not monotone game-by-game, and the exploration
noise of Chapter~\ref{ch:engineering} §8.5 is not optional --- without it the loop can
converge to a confident, narrow, and beatable policy. "Self-improvement" is a property of the
whole system at scale, not a magic property of the loss function.
\end{researchfinding}

\begin{pitfall}
A network trained on its own games is calibrated to \emph{its own} opponent distribution. That
is why the win probabilities of Chapter~\ref{ch:currency} are win probabilities \emph{between
two copies of Leela at the search depth used in training} --- not between two humans, and not
between you and your Saturday opponent. Any tool that shows a user "you had a $71\%$ win
probability here" is quietly changing the subject unless it says whose.
\end{pitfall}

\begin{repobox}[title={in this repository}]
Both heads are already extracted here: policy drives the arrows and the policy-blindness
metric; WDL drives the evaluation display and the "character of the fight". The moves-left
head appears to be unused --- it is free information, already computed on every evaluation,
and it is the natural signal for "is this position resolving or drifting?", which is a
question the drill scheduler asks in other ways.
\end{repobox}

\begin{exercises}

\exhead[the move vocabulary]{ch12:vocab} Why does the policy head have a fixed 1858 slots
rather than one output per legal move? Give a concrete reason connected to how a neural
network's output layer is built.

\exhead[masking]{ch12:masking} What would happen if illegal moves were \emph{not} masked
before the softmax --- to the probabilities, and to the search? Would the network learn
legality on its own? Should it have to?

\exhead[softmax over legal moves]{ch12:softmaxlegal} In a position with 46 legal moves, the
top prior is $31.93\%$. If the raw logit of that move were 1.0 higher, what would its
probability become, holding other logits fixed? (Hint: work with the ratio of exponentials;
assume the rest sums to $0.6807$ before renormalising.)

\exhead[disagreement]{ch12:disagree} Explain, in your own words and in three sentences, how
the same network gives Kd5 a $5.26\%$ prior while evaluating the position it reaches at
exactly $0.000$ with $d = 1.000$. Then say which of the two you would trust more, and why.

\exhead[moves left]{ch12:mlh} The moves-left head predicts $86.6$ plies in a position with
mate in two. Is this a failure of the head or a failure of something else? What would the
number become after search, and why?

\exhead[the training target]{ch12:target2} In the initial position, raw $P(\text{e4}) = 22.2\%$
and $\pi(\text{e4}) = 35.7\%$. Which is the training target for which? Write the direction of
the update in one sentence, and say what the network's policy for e4 would drift towards after
many such examples.

\exhead[why it is not circular]{ch12:notcircular} A sceptic says: "the network trains on its
own opinion, so it can only reinforce its errors." Answer in one paragraph, using the phrase
\emph{policy improvement operator} and referring to a specific number from
Chapter~\ref{ch:byhand}.

\exhead[the role of noise]{ch12:noiserole} Using Chapter~\ref{ch:engineering}'s
Equation~\eqref{eq:dirichlet}, explain what would go wrong in the training loop with no
Dirichlet noise at all. Connect your answer to the greedy failure of
Chapter~\ref{ch:manymachines} §4.2.

\exhead[whose chess?]{ch12:whosechess} Write two versions of a tooltip for the
policy-blindness metric: one that is precisely true, and one that is what most users would
assume it means. How far apart are they? In which kind of position do they diverge most?

\exhead[a use for M]{ch12:usem} Propose one concrete feature for your training tool that uses
the moves-left head, and one way that feature could mislead a user. Say how you would test it
before shipping.

\exhead[the loop, for humans]{ch12:humanloop} Equation~\eqref{eq:azloss} has two terms: match
the search's move preferences, and match the outcome. Describe the human training exercise
that corresponds to each term. Which of the two does your project's drill system currently do
better, and what would the other look like?

\end{exercises}
